% ======================================================================
% SEÇÃO 6 — CONCLUSÃO
% ======================================================================
\section{Conclus\~ao}
\label{sec:concl}

Este artigo apresentou o projeto, a implementa\c{c}\~ao e a valida\c{c}\~ao de um
sistema de automa\c{c}\~ao e supervis\~ao para a prepara\c{c}\~ao de amostras em
especia\c{c}\~ao de merc\'urio, concebido para substituir um sistema legado
baseado em LabVIEW. A solu\c{c}\~ao adota uma arquitetura mestre-escravo composta
por um Raspberry Pi (l\'ogica de alto n\'ivel em Python), um Arduino Uno
(aquisi\c{c}\~ao em tempo real) e uma IHM Web em React e TypeScript, integrados por
um protocolo JSON a 4~Hz.

O objetivo proposto foi atingido no \^ambito da valida\c{c}\~ao realizada: o
processo foi modelado em uma m\'aquina de estados finita com matriz de
acionamento; o controle PID misto da rampa do Tubo U e o controle do Forno 2 a
$700\,\celsius$ foram implementados e verificados por testes; o protocolo JSON, o
firmware do \daq (leitura SPI, atuadores e watchdog), a persist\^encia at\^omica
dos par\^ametros e a IHM Web foram desenvolvidos e integrados. A valida\c{c}\~ao
por 67 testes automatizados aprovados demonstrou a corre\c{c}\~ao da l\'ogica de
controle, a integra\c{c}\~ao entre as camadas e a efetividade dos mecanismos de
seguran\c{c}a no n\'ivel de software.

As principais contribui\c{c}\~oes s\~ao: a modelagem expl\'icita e verific\'avel do
sequenciamento; a estrat\'egia de controle mista que contorna a limita\c{c}\~ao de
medi\c{c}\~ao na regi\~ao criog\^enica; a arquitetura mestre-escravo de baixo
custo com separa\c{c}\~ao entre supervis\~ao e tempo real; o protocolo JSON aberto;
e o projeto de seguran\c{c}a em camadas (estado seguro por padr\~ao e watchdog
duplo).

Como limita\c{c}\~ao, a valida\c{c}\~ao foi conduzida em ambiente simulado, sem o
processo f\'isico; os crit\'erios de desempenho t\'ermico e a calibra\c{c}\~ao da
curva de aquecimento do Tubo U permanecem pendentes de verifica\c{c}\~ao
experimental. Como trabalhos futuros, sugere-se: validar o sistema em bancada;
integrar automaticamente o detector Lumex e exportar as curvas geradas;
implementar autentica\c{c}\~ao de operadores com trilha de auditoria;
desenvolver um simulador de processo para treinamento; e implementar a
sintoniza\c{c}\~ao autom\'atica dos controladores PID.
